\section{Open and closed subsets of spectra}
\label{section-open-and-closed}

\noindent
It turns out that open and closed subsets of a spectrum correspond to
idempotents of the ring.

\begin{lemma}
\label{lemma-idempotent-spec}
Let $R$ be a ring. Let $e \in R$ be an idempotent.
In this case
$$
\Spec(R) = D(e) \amalg D(1-e).
$$
\end{lemma}

\begin{proof}
Note that an idempotent $e$ of a domain is either $1$ or $0$.
Hence we see that
\begin{eqnarray*}
D(e)
& = &
\{ \mathfrak p \in \Spec(R)
\mid
e \not\in \mathfrak p \} \\
& = &
\{ \mathfrak p \in \Spec(R)
\mid
e \not = 0\text{ in }\kappa(\mathfrak p) \} \\
& = &
\{ \mathfrak p \in \Spec(R)
\mid
e = 1\text{ in }\kappa(\mathfrak p) \}
\end{eqnarray*}
Similarly we have
\begin{eqnarray*}
D(1-e)
& = &
\{ \mathfrak p \in \Spec(R)
\mid
1 - e \not\in \mathfrak p \} \\
& = &
\{ \mathfrak p \in \Spec(R)
\mid
e \not = 1\text{ in }\kappa(\mathfrak p) \} \\
& = &
\{ \mathfrak p \in \Spec(R)
\mid
e = 0\text{ in }\kappa(\mathfrak p) \}
\end{eqnarray*}
Since the image of $e$ in any residue field is either $1$ or $0$
we deduce that $D(e)$ and $D(1-e)$ cover all of $\Spec(R)$.
\end{proof}

\begin{lemma}
\label{lemma-spec-product}
Let $R_1$ and $R_2$ be rings.
Let $R = R_1 \times R_2$.
The maps $R \to R_1$, $(x, y) \mapsto x$ and $R \to R_2$,
$(x, y) \mapsto y$
induce continuous maps $\Spec(R_1) \to \Spec(R)$ and
$\Spec(R_2) \to \Spec(R)$.
The induced map
$$
\Spec(R_1) \amalg \Spec(R_2)
\longrightarrow
\Spec(R)
$$
is a homeomorphism. In other words,
the spectrum of $R = R_1\times R_2$ is the
disjoint union of the spectrum of $R_1$ and the
spectrum of $R_2$.
\end{lemma}

\begin{proof}
Write $1 = e_1 + e_2$ with $e_1 = (1, 0)$ and $e_2 = (0, 1)$.
Note that $e_1$ and $e_2 = 1 - e_1$ are idempotents.
We leave it to the reader to show that
$R_1 = R_{e_1}$ is the localization of $R$ at $e_1$.
Similarly for $e_2$.
Thus the statement of the lemma follows from Lemma
\ref{lemma-idempotent-spec} combined with Lemma
\ref{lemma-standard-open}.
\end{proof}

\noindent
We reprove the following lemma later after introducing
a glueing lemma for functions. See Section
\ref{section-tilde-module-sheaf}.

\begin{lemma}
\label{lemma-disjoint-decomposition}
Let $R$ be a ring. For each $U \subset \Spec(R)$
which is open and closed
there exists a unique idempotent $e \in R$ such that
$U = D(e)$. This induces a 1-1 correspondence between
open and closed subsets $U \subset \Spec(R)$ and
idempotents $e \in R$.
\end{lemma}

\begin{proof}
Let $U \subset \Spec(R)$ be open and closed.
Since $U$ is closed it is quasi-compact by
Lemma \ref{lemma-quasi-compact}, and similarly for
its complement.
Write $U = \bigcup_{i = 1}^n D(f_i)$ as a finite union of standard opens.
Similarly, write $\Spec(R) \setminus U = \bigcup_{j = 1}^m D(g_j)$
as a finite union of standard opens. Since $\emptyset =
D(f_i) \cap D(g_j) = D(f_i g_j)$ we see that $f_i g_j$ is
nilpotent by Lemma \ref{lemma-Zariski-topology}.
Let $I = (f_1, \ldots, f_n) \subset R$ and let
$J = (g_1, \ldots, g_m) \subset R$.
Note that $V(J)$ equals $U$, that $V(I)$
equals the complement of $U$, so $\Spec(R) = V(I) \amalg V(J)$.
By the remark on nilpotency above,
we see that $(IJ)^N = (0)$ for some sufficiently large integer $N$.
Since $\bigcup D(f_i) \cup \bigcup D(g_j) = \Spec(R)$
we see that $I + J = R$, see Lemma \ref{lemma-Zariski-topology}.
By raising this equation to the $2N$th power we conclude that
$I^N + J^N = R$. Write $1 = x + y$ with $x \in I^N$ and $y \in J^N$.
Then $0 = xy = x(1 - x)$ as $I^N J^N = (0)$. Thus $x = x^2$
is idempotent and contained
in $I^N \subset I$. The idempotent $y = 1 - x$ is contained in $J^N \subset J$. 
This shows that the idempotent $x$ maps to $1$ in every residue field
$\kappa(\mathfrak p)$ for $\mathfrak p \in V(J)$ and that $x$ maps to $0$
in $\kappa(\mathfrak p)$ for every $\mathfrak p \in V(I)$.

\medskip\noindent
To see uniqueness suppose that $e_1, e_2$ are
distinct idempotents in $R$. We have to show there
exists a prime $\mathfrak p$ such that $e_1 \in \mathfrak p$
and $e_2 \not \in \mathfrak p$, or conversely.
Write $e_i' = 1 - e_i$. If $e_1 \not = e_2$, then
$0 \not = e_1 - e_2  = e_1(e_2 + e_2') - (e_1 + e_1')e_2
= e_1 e_2' - e_1' e_2$. Hence either the idempotent
$e_1 e_2' \not = 0$ or $e_1' e_2 \not = 0$. An idempotent
is not nilpotent, and hence we find a prime
$\mathfrak p$ such that either $e_1e_2' \not \in \mathfrak p$
or $e_1'e_2 \not \in \mathfrak p$, by Lemma \ref{lemma-Zariski-topology}.
It is easy to see this gives the desired prime.
\end{proof}

\begin{lemma}
\label{lemma-characterize-spec-connected}
Let $R$ be a nonzero ring. Then $\Spec(R)$ is
connected if and only if $R$ has no nontrivial
idempotents.
\end{lemma}

\begin{proof}
Obvious from Lemma \ref{lemma-disjoint-decomposition}
and the definition of a connected topological space.
\end{proof}

\begin{lemma}
\label{lemma-ideal-is-squared-union-connected}
Let $I \subset R$ be a finitely generated ideal of a ring $R$
such that $I = I^2$. Then
\begin{enumerate}
\item there exists an idempotent $e \in R$ such that $I = (e)$,
\item $R/I \cong R_{e'}$ for the idempotent $e' = 1 - e \in R$, and
\item $V(I)$ is open and closed in $\Spec(R)$.
\end{enumerate}
\end{lemma}

\begin{proof}
By Nakayama's Lemma \ref{lemma-NAK} there exists an element
$f = 1 + i$, $i \in I$ such that $fI = 0$. Then $f^2 = f + fi = f$
is an idempotent. Consider the idempotent $e = 1 - f = -i \in I$.
For $j \in I$ we have $ej = j - fj = j$ hence $I = (e)$.
This proves (1).

\medskip\noindent
Parts (2) and (3) follow from (1). Namely, we have
$V(I) = V(e) = \Spec(R) \setminus D(e)$ which is open and
closed by either Lemma \ref{lemma-idempotent-spec} or
Lemma \ref{lemma-disjoint-decomposition}. This proves (3).
For (2) observe that the map $R \to R_{e'}$ is surjective
since $x/(e')^n = x/e' = xe'/(e')^2 = xe'/e' = x/1$ in $R_{e'}$.
The kernel of the map $R \to R_{e'}$ is the set of elements of
$R$ annihilated by a positive power of $e'$. Since $e'$ is
idempotent this is the ideal of elements annihilated by $e'$
which is the ideal $I = (e)$ as $e + e' = 1$ is a pair
of orthogonal idempotents. This proves (2).
\end{proof}






